
\section*{Reviewer 2}
学習アルゴリズムの修正については R1.1 と R1.24、以下で問題とされた終端規則の完全性の補修については R2.8 を参照されたい。

\subsection*{固定 $h$ の全般的な位置づけ}
\Comment{R2.1}
「大局的には本論文は、固定長窓を任意の認識可能な合同に置き換えることで Yoshinaka の $(k,\ell)$-代入可能言語の学習を抽象化し、既知の結果を一般化している。この抽象化は強力で十分に理解された枠組みに問題を置く一方、先行研究がクラス全体に一様なパラメータを固定したのに対し、各言語に対して準同型を丸ごと事前に知ることを要求する。準同型の学習が範囲外だと認めているが、$(k,\ell)$-代入可能言語の研究のための抽象化以上に、なぜこの設定が興味深いのか明らかでない。有望な方向は、第8--10節でも触れているプッシュダウン・オートマトンとの関係を深めることだろう。Takada の研究との関係も調べる価値がある。……より実用的な方向では、本研究をオートマトンの型付けを文脈自由代入可能言語へ拡張するものとみなし、型付けが満たすべき条件を調べることもできる。」

\Reply
有限モノイド準同型の指定は $(k,\ell)$ の選択より強い事前バイアスであり，学習器は $h$ 自体を推定しない。ただし一つの $h$ は $\Chcf$ 全体に対して一様に固定される。命題8.4は固定窓のどのクラスにも属さない非正規線形例を示し，系7.5と定理8.3は窓・有限群型付けおよび線形の場合に符号長で多項式の特徴データを保証する。型付けを動かす合併 $\RS$ の学習可能性は主張しない。$h=h_{k,\ell}$ のときの再構成は Yoshinaka (2008, Section 4) の部分語規則方式を特殊化する（表現上の差異は第4.1節）。

序論には Kanazawa (1996; 1998, Chapters 5, 8) の範疇文法の学習を相補的な先行研究として挙げたが，それらから本稿の固定 $h$ 定理が従うとは主張しない。Yoshinaka (2015, Theorem 1) の一般的多ソート枠組みは，固定 $h$ の非空繊維 $h^{-1}(m)\cap\Sigma^+$ をソートとすることで定性的学習可能性を包含する。この型の選択をYoshinaka本人に帰属させる趣旨ではない。有限モノイドに固有の定量的・構造的結果とは区別する。

第3節では Yoshinaka (2008, Proposition 1) の非閉包例と命題3.2を対比し，消去を伴う逆像に対しては型付けを変更して閉包性を得るのであって，元の固定窓や固定 $h$ のクラスの閉包性を主張しないことを明示した。命題3.2(iii)の脚注は消去フラグ $e_\varphi$ の必要性を例示する。局所自明半群の判定条件は命題3.5でPinに帰属させ，補題7.4は局所自明半群と群の直積へのリフトによる窓・有限群の核の特徴づけを与える。系7.5の通常厚みに関する上界は型付き生成語の評価と正規化を利用する。命題6.3・補題7.6・系6.4は，二元半束型付けによる指数的特徴データと型付き厚みの下界を与えるが，任意の型付けについて代数的条件の必要性やGold学習器一般の実行時間・収束時刻下界は主張しない。

序論には Coste, Fredouille, Kermorvant, de la Higuera (2004) の自動機の型付けを引用した。型付けの適切さとは分布上の衝突を有限指数のモノイド合同で分離することであり，構文合同そのものを精密化する条件は正規性を強制する。これは型付けを見つけるアルゴリズムではない。第8節のTakadaの正規制御集合との関係は限定的であり，一般のプッシュダウン特徴づけや同値性は主張しない。
\Location{序論、第3節（命題3.2直前の Yoshinaka との比較）、第4節、命題3.2・3.5、主定理3.4(iv)、命題6.3・系6.4・付録C、第7節・補題7.4・7.6・系7.5・付録A.1、命題8.4、第9.2節（$L_{\mathrm{Yreg}}$ の定義）、第8節の制御集合に関する議論。}

\subsection*{主要な懸念1：外側文脈の型付けと輸送}
\Comment{R2.2}
「最大の懸念は型付き精密化、とくに輸送にある。規則(2)は外側の型に何の制約もないまま、非終端記号を任意の文脈間で移動させる。これは意図されたことなのか。外側の成分 $m,n$ が非終端記号の出現位置を記録するなら、それを移動させる規則が，なぜその成分を完全に無視するのか。私の見るところ、この成分は証明で一度も参照されず、すべての $A^{m,n}_p$ を $A_p$ に置き換えても結果は変わらない。著者は同意するか。そうでなければ、外側の型がどこで参照されるのか示してほしい。これは定義2.3そのものから従うように思える。代数的条件が課されるのは置換される因子だけであり、その周辺文脈ではない。ゆえにその定義を模倣する学習器では、輸送を制限する条件は何もないはずだ。もしそうなら、第1.1節で予告する『型付き輸送』は論文のどこにも登場せず、出現を外側の型に応じて分離すべきだと示す例も存在しえない。自明なモノイドで空虚となる生成語型付けは本当に必要なので、それを脚注のままにしておくべきではない。」

\Reply
同意する。未使用の外側文脈型注釈を削除し，標的側の $A_p^{m,n}$ を生成語型付きの $A_p$ に単純化し，仮説側の $[x:u,v]$ を観測部分語 $x$ ごとの $\NT{x}$ にまとめた。単項規則は標本の同一文脈と $h^\sharp$-型一致を使い（非空因子では元の $h$-型一致と同値），独立の文脈輸送規則は不要となった。投稿時の外側文脈型の必要性の主張を撤回した。健全性は定理4.2の不変量 $\NT{x}\Rightarrow^*w$ ならば $x\equiv_Lw$ かつ $h^\sharp(w)=h^\sharp(x)$ により証明し，外側の型を参照しない。

一方，生成語型付けは本稿の有限証拠集合に基づく完全性証明に必要である。第5.1節に本文で説明し，標準親生成語から型付き終端生成語または子生成語の連結へ進む単項規則(U)を許す。偶奇型付けの例はこの証明での必要性を説明するのであって，任意のアルゴリズムへの普遍的な必要性を主張するものではない。
\Location{第4節、第5.1節、定理5.3。}

\subsection*{主要な懸念2：線形定理の射程}
\Comment{R2.3}
「線形の場合の結果には別の疑問がある。対象のクラスはどれほど大きいのか。末尾の三つの節は $\RS$ の位置づけを述べるだけで、$\mathcal C_h^{\mathrm{lin}}$ について述べておらず、定理7.9が通常の線形代入可能言語の先にどこまで届くのか読者にはわからない。上限付きカウンター言語は正規だから線形でもあり、Yoshinaka の階層との分離はすでに線形クラスの中で成り立つ。一方、$\{a^nb^n\}^*$ は線形ではないため、最後の節は線形のレベルについて何も示していない。現状では、非正規な線形の場合を未解決に残し、射程不明なクラスに対する技術的な強化に見える。」

\Reply
第8節に具体的な非正規線形例と正規フィルターによる構成を加えた。命題8.4の $L_{\pm,e}$ は固定有限モノイド型付けで代入可能な線形・非正規言語だが，Clark--Eyraudクラスにも任意の $(k,\ell)$-クラスにも属さない。命題3.2(ii)と線形言語の正規言語との共通部分の閉包性により，$L\in\Chlin$ かつ $R_{\mathrm{reg}}=g^{-1}(\mathsf{Acc}_g)$ ならば $L\cap R_{\mathrm{reg}}\in\mathcal C_{h\times g}^{\mathrm{lin}}$ である。この例の証明は $L_{\pm,e}=L_{\mathrm{all}}\cap R_{\mathrm{reg}}$ と書き，別に四元モノイドの具体的な型付けも示す。正規言語段階の分離は Yoshinaka (2008, Proposition 1) の既知結果に帰する。

定理8.3は固定 $h$ の代入可能な線形言語に対し，選んだ線形CFGに相対的な集合依存 $\Bh$ の多項式特徴データを保証する。これは一般の線形言語についてのde la Higueraの否定的結果に反しない。$\Ah$ の逐次的な多項式データ保証も主張しない。
\Location{命題3.2(ii)、定理8.3とその適用範囲の補足、命題8.4（証明中の正規フィルターの議論）および後続の四元型付けの段落。}

\subsection*{その他の指摘}
\Comment{R2.4}
「この論文はもっと簡潔にすると改善される。第4節の多くは定義であり、いくつかの結果は直前の証明で示したことを繰り返している。対象を導入する前に使用することも多く、読者は未解決のまま説明を読み進めなければならない。また $\widetilde G$ は $\widetilde G^{\mathrm{full}}$ の刈り込みであるから、記号が到達可能かつ生成可能なのは定義上当然であり、毎回断る必要はなく、その規約を一文で述べれば十分だ。」

\Reply
同意する。到達可能性・生成可能性の繰り返しの限定を削除し，第5.1節で縮約済み精密化の規約を一度だけ述べた。旧補題4.6--4.7，5.3--5.4および定理4.9，4.12などの冗長な結果を削除した（R1.22，R2.9）。再構成作用素は標的側精密化より先に，$\omega$ は初出前に，$\chi$ は証拠集合の直前に，\emph{locally trivial} は命題3.5より先の第2.2節で定義した。定量的議論を第6節（一般の型付き厚み），第7節（窓・有限群），第8節（線形言語）に分け，長い正規化の証明は付録に移した。
\Location{第4節、第5.1節、第6--8節、付録A--B。}

\Comment{R2.5}
「命題3.1：$L=h^{-1}(F)$ という仮定は正規言語を正確に特徴づける。『In particular』ではなく『Thus』とすべきで、証明で理由を明示してもよい。」

\Reply
同意する。命題3.1は有限モノイドによる認識可能性の特徴づけを直接述べ、そこからすべての正規言語が $\RS$ に属すると結論する。標準的な特徴づけについて Eilenberg (1974, Chapter III, Propositions 10.1--10.2) を引用し、必要な帰結も示す。$L=h^{-1}(\mathrm{Acc})$ で $h(x)=h(y)$ なら、すべての文脈 $(u,v)$ に対して $h(uxv)=h(uyv)$ となるから分布は等しい。
\Location{命題3.1。}

\Comment{R2.6}
「定義4.11：基底 $\mathsf B(\widetilde G)$ は使われていない。定理4.12、注4.13とともに導入する必要があるのか。$\mathrm{NT}(\cdot)$ と $\mathrm{Rule}(\cdot)$ も不要であり、$G$ を組として定義し、その $V,P$ を使えば十分だ。」

\Reply
同意する。これらの対象と記法は削除した。
\Location{第5.1節。}

\Comment{R2.7}
「例4.3は SSBNF ではない。」

\Reply
旧例は全面的に削除した。改訂後の証明では不要となったためである。
\Location{第5.1節。}

\Comment{R2.8}
「補題5.2(i)：$X$ に複数の終端規則があるなら $\omega(X)$ だけでは不十分で、規則(3)を使うべきだ。証明を修正する必要がある。」

\Reply
指摘は正しい。投稿時の完全性証明には実際に欠落があった。型付き終端規則 $X\to a$ ごとに同じ標本文脈の証拠語 $u_Xav_X$ を明示的に含め，標準証拠語 $u_X\omega(X)v_X$ と合わせて，必要に応じて単項規則(U)を $\NT{\omega(X)}$ から $\NT{a}$ に適用し，語彙規則(L)を使う。型付き精緻化による $h^\sharp(\omega(X))=h^\sharp(a)$ がその単項規則を正当化する。
\Location{第5.1節（証拠集合）、定理5.3（終端規則の基本場合）。}

\Comment{R2.9}
「補題5.3と5.4は必要か。」

\Reply
同意する。両補題は削除した。
\Location{第4--5.1節。}

\Comment{R2.10}
「系5.8は『有限モノイド $M$ への固定準同型 $h$ について……』、または『準同型 $h$ が与えられているとき……』という文で始めるべきだ。」

\Reply
改訂後の Gold 同定の文は、提案された『有限モノイド $M$ への固定準同型 $h:\Sigma^*\to M$ について……』という形で始める。
\Location{系5.5。}

\Comment{R2.11}
「第7.1節の標準 SSLNF への変換（3ページ）は、短縮するか付録に移せるだろう。」

\Reply
付録に移す案を採用した。第8節には定理8.3に直接必要な命題8.1・補題8.2を残し，補題8.2の短い証明も本文に置いた。正規化の証明は付録Bにまとめた。
\Location{命題8.1と周辺、補題8.2とその証明、定理8.3とその後の補足、付録B。}
